\section{Background}

\subsection{Value Prediction}

Value prediction estimates an instruction's result before it becomes
available, potentially reducing the latency experienced by dependent
instructions.

Predicted values require explicit rules for their use and validation.
When an incorrect prediction has been used, the processor must also
preserve correctness through appropriate recovery. Validation may
check the predicted value directly or, as in V-ASTRA, validate the
address from which the value was obtained.

\subsection{Last Value Prediction}

Last Value Prediction (LVP) predicts an instruction's result using
the value produced by its most recent instance. It is a computational
predictor that applies the identity function to the last observed
value~\cite{ravenel2025architecting}. Predictions are typically
associated with static instructions using their PCs, allowing
observations from earlier executions to be reused.

Since instruction values may change across executions, LVP uses a
confidence counter to control prediction admission. The design in
\cite{ravenel2025architecting} uses a two-bit saturating counter,
incremented on correct predictions and decremented on incorrect
ones. A prediction is admitted only when the counter is saturated.
In this implementation, the admission threshold is configurable.
The counter is updated for every eligible load that completes,
regardless of whether a prediction was made or used. Thus, entries
below the threshold can continue to gain confidence.

Prediction accuracy and coverage are both important, but neither
alone determines performance. High accuracy with few predictions
may provide limited benefit, while high coverage may be ineffective
if the predicted instructions do not lie on critical dependency
chains. The resulting cycle-count improvement must therefore also
be evaluated.

LVP is particularly relevant to in-order processors, where a
long-latency load can delay younger dependent instructions. Using
a predicted value allows eligible consumers, including dependent
memory operations, to proceed before the producer completes,
potentially overlapping memory operations that would otherwise
wait for one another.

A prediction does not replace the actual result. When the producer
completes, its actual value is compared with the prediction. If an
incorrect value was used, the processor must recover the affected
younger instructions to preserve architectural correctness.

\subsection{ASTRA Address Prediction}

ASTRA (Address Stride Alias) is the address predictor proposed
in~\cite{ravenel2025architecting}. It uses an order-1 differential
finite-context method (D-FCM), with two modifications: it uses the
last observed stride directly as history rather than hashing past
strides, and its tables are untagged.

ASTRA is indexed by load PC. Its Last Address Table (LAT) stores
the last committed virtual address and signed stride for each
entry. The stride indexes a second table, Stride$\rightarrow$Stride
(SS), which provides the stride added to the LAT address to form
the predicted address. On load retirement, ASTRA computes the new
stride, updates the corresponding SS entry, and records the
committed address and stride in the LAT.

The untagged SS table allows loads observing the same stride to
share entries. The paper exploits this constructive aliasing so
that an update by one load can help other loads predict a new
access pattern. ASTRA therefore favors prediction coverage over
accuracy, relying on address validation before a prediction is
used (Section~\ref{sec:bg-vastra}).

The original ASTRA predicts the complete load address, whereas
V-ASTRA predicts the load's base-register value and adds the
instruction displacement. This project implements the
base-address form: its LAT stores the base-register value rather
than the complete effective address
(Section~\ref{sec:impl-astra}). ASTRA predicts an address, not
the loaded data. Obtaining the load value requires V-ASTRA's
separate value cache.

\subsection{V-ASTRA and the Value Cache}
\label{sec:bg-vastra}

V-ASTRA (Value-Address Stride Alias) uses ASTRA's predicted address
to access a dedicated value cache, which stores previously
observed load values. In the paper, the value cache holds 4\,KiB of data
in 16-byte entries. Loads allocate entries at commit, while stores
update matching entries without allocating new ones.

V-ASTRA predicts a load's base address and adds the displacement
from the decoded instruction to obtain its effective address.
The paper validates the predicted base at issue by comparing it
with the actual base-register value. An incorrect prediction is
therefore rejected before use, requiring neither a confidence
counter nor value-misprediction recovery.

Safe bypass also requires the value cache to contain the requested
bytes and no older in-flight store to modify them. The paper uses
a counting Bloom filter to detect potentially conflicting stores
and an invalidation FIFO to handle older stores issuing after a
prediction. When all checks succeed, the load skips the L1D access
and its dependents can issue in the same cycle, providing
zero-cycle load-use behavior~\cite{ravenel2025architecting}.

The implementation in this work differs from the paper. Candidate
selection occurs in Execute when the load computes its address,
while actual bypass acceptance occurs later in the LSQ, before
the normal memory request is sent. It also uses an exact scan of
older stores instead of the paper's counting Bloom filter and
invalidation FIFO (Section~\ref{sec:impl-vcache}).

\subsection{Dependency Handling}

A producer-consumer dependency occurs when an instruction requires
a value produced by an earlier instruction. In an in-order processor,
a consumer waiting for a long-latency load can also delay younger
independent instructions because issue follows program order. Under
stall-on-use, a dependent instruction cannot issue until the
producer's value becomes available.

Two mechanisms can reduce this delay. Dependency release
makes the consumer eligible to issue earlier but does not provide
the required value. Value forwarding supplies the producer's
value directly, allowing an eligible consumer to proceed before the
producer completes. While the paper's issue-time V-ASTRA validation
allows dependents to issue in the same cycle as the load, MinorCPU
handles operands differently: ordinary non-memory instructions read
their operands at commit, whereas memory operations can issue to
the LSQ earlier (Section~\ref{sec:bg-minor}). Consequently, the
forwarding mechanism in this work is restricted to eligible
dependent loads.

\subsection{META Selection}

ASTRA/V-ASTRA and LVP exploit different forms of predictability:
memory-access patterns and repeated instruction values, respectively.
META combines them to improve prediction coverage
\cite{ravenel2025architecting}.

In the original design, both predictors operate in parallel, and
META selects a prediction at issue time. A validated V-ASTRA value
is preferred when the predicted base is correct and the requested
data is available in the value cache. Otherwise, META uses an
LVP prediction if its confidence counter is saturated. If neither
is available, the load is not predicted. Both predictors are updated
when the load retires~\cite{ravenel2025architecting}.

The two sources differ in their validation requirements. V-ASTRA
validates its address before use, whereas LVP remains speculative
until the actual value is available. An incorrect LVP prediction
therefore requires recovery only if it was used. The paper replays
younger instructions through the scoreboard, while this work
squashes and refetches them.

This project adapts META to MinorCPU by selecting the source of
a dependent load's base operand when it attempts to proceed before
its producer commits. A validated V-ASTRA value is preferred.
If the V-ASTRA outcome is Pending, the consumer waits rather than
falling back to LVP. If V-ASTRA has Failed or the producer is not
a candidate, an eligible LVP prediction may be used. Otherwise,
the consumer follows normal execution. This pending-state rule is
specific to the implementation: the paper validates V-ASTRA at
issue, so this undecided state does not arise.

\subsection{gem5 MinorCPU}
\label{sec:bg-minor}

The implementation uses MinorCPU, gem5's in-order processor model
\cite{gem5}. Its four pipeline stages are Fetch1 (instruction fetch),
Fetch2 (instruction separation and branch prediction), Decode, and
Execute. Execute contains the functional-unit pipelines, in-flight
instruction queue, commit logic, and load-store queue (LSQ).

Instructions issue to functional units in program order when their
source registers are ready according to the scoreboard. The
scoreboard tracks expected result availability and the latest
in-flight writer of each register, but does not hold register values.
Results are written to architectural registers at commit, also in
program order.

Memory instructions can compute their addresses and enter the LSQ
before reaching the head of the in-flight queue, but issue one at a
time and ordinarily only after their source-producing instructions
have committed. Loads write their destination registers at commit.
Branches are resolved in Execute, with mispredictions causing
younger wrong-path instructions to be discarded.

The memory functional unit uses \texttt{extraAssumedLat} to estimate
when dependents may issue. With the default setting, dependents may
issue after an assumed load latency, but wait for the producer to
commit before executing. Setting this latency to zero enables the
stall-on-use configuration, in which dependents wait for commit
before issuing.

These properties constrain prediction integration. This work adds
LVP operand substitution with squash-based recovery and forwarding
of validated V-ASTRA values to dependent loads. Unlike the paper,
which validates the predicted base at issue, this implementation
compares it with the actual base when Execute initiates the memory
access. Bypass acceptance occurs later in the LSQ, before the normal
memory request is sent. Address translation still precedes bypass.
These differences make the implementation an adaptation of the
paper's mechanisms rather than a reproduction of its complete
processor architecture and evaluation environment.
